% !PRINTED 336
\weberpage{336}
If, then, we denote variables by $x_r$ and set
\weq{7}{\overset{r}{\Sigma}\, x_r\, c^{(r)}_{h,k} = y_{h,k},}
it follows that:
\weq{8}{x_s = \overset{h,k}{\Sigma}\, a^{(s)}_{h,k}\, y_{h,k}.}
Now let $u_r$, $v_r$, $t_r$ denote three systems of variables, and let
\weq{9}{\begin{aligned}
\mu &= \Sigma\, \alpha_r\, u_r, \\
\nu &= \Sigma\, \beta_r\, v_r, \\
\lambda &= \Sigma\, \gamma_r\, t_r
\end{aligned}}
be basis forms of $\mathfrak{a}$, $\mathfrak{b}$, $\mathfrak{c}$, and further let $U$, $V$, $T$ be the forms belonging to $\mathfrak{a}$, $\mathfrak{b}$, $\mathfrak{c}$, written in the variables $u,v,t$. Then
\weq{10}{\begin{aligned}
N(\mu) &= N(\mathfrak{a})\,U, \\
N(\nu) &= N(\mathfrak{b})\,V, \\
N(\lambda) &= N(\mathfrak{c})\,T.
\end{aligned}}
From (3) there follows therefore, by multiplication with $u_h v_k$ and summation according to (9):
\weq{11}{\mu\nu = \overset{r}{\Sigma}\, \gamma_r \overset{h,k}{\Sigma}\, c^{(r)}_{h,k}\, u_h v_k,}
hence, if one sets
\weq{12}{t_r = \overset{h,k}{\Sigma}\, c^{(r)}_{h,k}\, u_h v_k}
then:
\weq{13}{\mu\nu = \lambda,}
and from this, since $N(\mathfrak{a})N(\mathfrak{b}) = N(\mathfrak{c})$, there follows by (10)
\weq{14}{T = UV.}
In this, however, the $t_r$ are no longer independent variables; rather, they arise from $u_h$, $v_k$ through the \gesperrt{bilinear substitution} (12).

5. \gesperrt{The form $T$ passes, through the bilinear substitution (12), into the product of the two forms $U, V$.}

But this substitution has still one essential peculiarity. We call the functions $t_r$ determined by (12) \gesperrt{linearly independent} with respect to a prime modulus $p$ if from the congruence
\weq{15}{\Sigma\, x_r t_r \equiv 0 \pmod{p},}
in which the $x_r$ are rational integers, it follows that all these numbers are divisible by $p$. If this holds for every arbitrary prime $p$, the $t_r$ are called \gesperrt{simply linearly independent}.
